\documentclass[12pt,a4paper]{article}

% --------------------------------------------------
% PACOTES BÁSICOS
% --------------------------------------------------
\usepackage[utf8]{inputenc}
\usepackage[T1]{fontenc}
\usepackage[brazil]{babel}
\usepackage{lmodern}
\usepackage[a4paper,margin=2.5cm]{geometry}
\usepackage{setspace}
\usepackage{titlesec}
\usepackage{tocloft}
\usepackage{array}
\usepackage{longtable}
\usepackage{booktabs}
\usepackage{tabularx}
\usepackage{enumitem}
\usepackage{hyperref}
\usepackage{xcolor}
\usepackage{fancyhdr}
\usepackage{lastpage}
\usepackage[backend=biber,style=ieee]{biblatex}
\addbibresource{sections/references.bib}
\input{config.tex}

% --------------------------------------------------
% CONFIGURAÇÕES GERAIS
% --------------------------------------------------
\onehalfspacing
\setlength{\parindent}{1.25cm}
\setlength{\parskip}{0.2cm}

\hypersetup{
    colorlinks=true,
    linkcolor=black,
    urlcolor=blue,
    citecolor=black,
    pdftitle={Especificação dos Requisitos},
    pdfauthor={Autores do Projeto}
}

% --------------------------------------------------
% CABEÇALHO E RODAPÉ
% --------------------------------------------------
\pagestyle{fancy}
\fancyhf{}
\fancyhead[L]{Especificação dos Requisitos do \ProjetoNome}
\fancyhead[R]{Página \thepage}
\renewcommand{\headrulewidth}{0.4pt}

% --------------------------------------------------
% FORMATAÇÃO DE TÍTULOS
% --------------------------------------------------
\titleformat{\section}
  {\normalfont\bfseries\large}
  {\thesection}{1em}{}

\titleformat{\subsection}
  {\normalfont\bfseries\normalsize}
  {\thesubsection}{1em}{}

\titleformat{\subsubsection}
  {\normalfont\bfseries\normalsize}
  {\thesubsubsection}{1em}{}

\newcommand{\blocoRF}[5]{
\noindent\textbf{#1 \textless #2\textgreater}

\vspace{0.2cm}
\noindent\textbf{Descrição:} #3

\noindent\textbf{Prioridade:} #4

\noindent\textbf{Ator(es):} #5

\noindent\textbf{Requisitos associados:} \textless Ex.: RF002, RF015 \textgreater

\vspace{0.4cm}
}

\newcommand{\blocoNF}[3]{
\noindent\textbf{#1 \textless #2\textgreater}

\vspace{0.2cm}
\noindent\textbf{Descrição:} #3

\vspace{0.4cm}
}

% --------------------------------------------------
% INÍCIO DO DOCUMENTO
% --------------------------------------------------
\begin{document}

% --------------------------------------------------
% CAPA
% --------------------------------------------------
\begin{titlepage}

\noindent\rule{\textwidth}{0.8pt} % traço no topo

\begin{flushright}
    {\Huge \textbf{Especificação dos Requisitos}}
\end{flushright}

\begin{flushright}
    {\Large \textbf{do}}\\[1.0cm]
    
    {\Huge \textbf{\ProjetoNome}}\\[1.2cm]
    
    {\Large \textbf{Versão \VersaoDocumento}}
\end{flushright}

\vspace{1.8cm}
\begin{flushleft}
    \AutorA\\
    \AutorB\\
    \AutorC\\[0.5cm]

    \textbf{Professoras:} \\
        \ProfessoraA\\
        \ProfessoraB\\[0.5cm]
    \textbf{Disciplina:} \DisciplinaNome
\end{flushleft}

\end{titlepage}

% --------------------------------------------------
% SUMÁRIO
% --------------------------------------------------
\tableofcontents
\newpage




% --------------------------------------------------
% REVISÕES
% --------------------------------------------------
\section*{Revisões}
\begin{longtable}{|p{2.5cm}|p{4cm}|p{5cm}|p{3cm}|}
\hline
\textbf{Versão} & \textbf{Autores} & \textbf{Descrição da Versão} & \textbf{Data} \\
\hline
1.0 & Equipe inicial & Criação inicial do documento & \textless dd/mm/aaaa\textgreater \\
\hline
1.1 & Equipe atualizada & Ajustes de conteúdo e refinamento dos requisitos & \textless dd/mm/aaaa\textgreater \\
\hline
\end{longtable}

\newpage

% --------------------------------------------------
% 1 INTRODUÇÃO
% --------------------------------------------------
\section{Introdução}
\input{sections/introduction}

\subsection{Objetivo do Documento}
\input{sections/purpose}

\subsection{Escopo do Produto}
\input{sections/scope}

\subsection{Público-Alvo}
\input{sections/audience}

\subsection{Definições, Acrônimos e Abreviações}
\input{sections/definitions}

\subsection{Convenções}
\input{sections/conventions}

\subsection{Referências}
\nocite{*}
\printbibliography[heading=none]




\newpage
% --------------------------------------------------
% 2 VISÃO GERAL
% --------------------------------------------------
\section{Visão Geral}

\subsection{Perspectiva do Produto}
Explicar a motivação do desenvolvimento do produto e descrever como ele irá interagir com o ambiente, sistemas externos, usuários e demais componentes envolvidos.

\subsection{Funcionalidade do Produto}
Descrever as principais funcionalidades do sistema em alto nível.

\begin{itemize}[leftmargin=1.2cm]
    \item Cadastro e autenticação de usuários
    \item Gerenciamento de dados do domínio do sistema
    \item Consulta e atualização de informações
    \item Emissão de relatórios ou visualizações
    \item Comunicação com serviços externos, se aplicável
\end{itemize}

\subsection{Usuários}
Descrever os perfis de usuários que irão utilizar o sistema.

\begin{itemize}[leftmargin=1.2cm]
    \item \textbf{Administrador}: gerencia usuários e configurações
    \item \textbf{Usuário comum}: utiliza as funcionalidades principais
    \item \textbf{Operador/Professor/Técnico}: perfil especializado, se houver
\end{itemize}

\subsection{Ambiente Operacional}
Descrever a plataforma, sistema operacional, hardware e demais aplicações que coexistirão com o produto.

\begin{itemize}[leftmargin=1.2cm]
    \item Sistema operacional: Windows, Linux, Android, iOS ou Web
    \item Linguagens e frameworks: conforme projeto
    \item Banco de dados: conforme projeto
    \item Navegadores compatíveis: Chrome, Firefox, Edge
\end{itemize}

\subsection{Restrições de Projeto e Implementação}
Descrever limitações técnicas e de projeto.

\begin{itemize}[leftmargin=1.2cm]
    \item Restrições de hardware
    \item Linguagem de programação definida
    \item Padrões arquiteturais
    \item Protocolos de comunicação
    \item Prazo acadêmico de entrega
\end{itemize}

\subsection{Documentação do Usuário}
Listar os documentos de apoio ao usuário.

\begin{itemize}[leftmargin=1.2cm]
    \item Manual do usuário
    \item Guia de instalação
    \item Guia rápido de uso
    \item Ajuda embutida no sistema
\end{itemize}

\subsection{Suposições e Dependências}
Descrever as suposições e dependências que podem afetar os requisitos.

\begin{itemize}[leftmargin=1.2cm]
    \item Disponibilidade de internet
    \item Acesso a banco de dados
    \item Dependência de APIs externas
    \item Disponibilidade dos usuários para validação
\end{itemize}

% --------------------------------------------------
% 3 ESPECIFICAÇÃO DAS INTERFACES EXTERNAS
% --------------------------------------------------
\section{Especificação das Interfaces Externas}

\subsection{Requisitos de Interface Externa}

\subsubsection{Interfaces do Usuário}
As interfaces do usuário devem ser intuitivas, responsivas e coerentes com o perfil dos usuários do sistema. Protótipos, wireframes ou telas podem ser entregues em anexo.

\subsubsection{Interfaces de Hardware}
Descrever as conexões do produto com hardwares utilizados, quando aplicável, como leitores, sensores, impressoras, dispositivos móveis ou periféricos específicos.

\subsubsection{Interfaces de Software}
Descrever as conexões com softwares externos, bibliotecas, frameworks, bancos de dados, serviços em nuvem ou APIs consumidas pelo sistema.

\subsubsection{Interfaces de Comunicação}
Descrever como serão realizadas as comunicações, como navegador web, e-mail, formulários eletrônicos, REST API, sockets ou protocolos específicos.

% --------------------------------------------------
% 4 REQUISITOS FUNCIONAIS
% --------------------------------------------------
\section{Requisitos Funcionais}

Nesta seção são definidas as características que o sistema deve apresentar por meio da especificação dos requisitos. Cada requisito deve possuir um identificador único. A nomenclatura recomendada é:

\begin{itemize}[leftmargin=1.2cm]
    \item Para requisitos funcionais: \textbf{[RF001] Nome do requisito}
    \item Para requisitos não-funcionais: \textbf{[NF001] Nome do requisito}
\end{itemize}

Para estabelecer a prioridade dos requisitos, adotam-se as denominações:

\begin{itemize}[leftmargin=1.2cm]
    \item \textbf{Essencial}: sem este requisito o sistema não entra em funcionamento de forma viável.
    \item \textbf{Importante}: sem este requisito o sistema funciona, mas de forma insatisfatória.
    \item \textbf{Desejável}: requisito útil, porém não indispensável para a versão atual.
\end{itemize}

\subsection{RF001 \textless Nome do requisito\textgreater}

\blocoRF
{RF001}
{Nome do requisito}
{o que representa no sistema e/ou o que deve fazer}
{Essencial}
{Informar os atores que irão interagir com o requisito}

\subsection{RF002 \textless Nome do requisito\textgreater}

\blocoRF
{RF002}
{Nome do requisito}
{o que representa no sistema e/ou o que deve fazer}
{Importante}
{Informar os atores que irão interagir com o requisito}

\subsection{RF003 \textless Nome do requisito\textgreater}

\blocoRF
{RF003}
{Nome do requisito}
{o que representa no sistema e/ou o que deve fazer}
{Desejável}
{Informar os atores que irão interagir com o requisito}

% --------------------------------------------------
% 5 REQUISITOS NÃO-FUNCIONAIS
% --------------------------------------------------
\section{Requisitos Não-Funcionais}

\subsection{Requisitos de Desempenho}

\blocoNF
{NF001}
{Tempo de resposta}
{O sistema deve responder às principais operações do usuário em até \textless X\textgreater\ segundos, considerando condições normais de uso.}

\blocoNF
{NF002}
{Capacidade}
{O sistema deve suportar simultaneamente \textless X\textgreater\ usuários conectados sem degradação crítica do serviço.}

\subsection{Requisitos de Segurança}

\blocoNF
{NF003}
{Controle de acesso}
{O sistema deve restringir o acesso às funcionalidades de acordo com o perfil do usuário autenticado.}

\blocoNF
{NF004}
{Proteção de dados}
{Os dados sensíveis devem ser armazenados e transmitidos de forma segura, utilizando mecanismos adequados de proteção.}

\subsection{Atributos de Qualidade do Software}

\blocoNF
{NF005}
{Usabilidade}
{O sistema deve apresentar interface clara, padronizada e de fácil aprendizado para o usuário.}

\blocoNF
{NF006}
{Manutenibilidade}
{O código-fonte deve ser modular, documentado e organizado para facilitar manutenção e evolução.}

\blocoNF
{NF007}
{Confiabilidade}
{O sistema deve manter consistência dos dados e se recuperar adequadamente de falhas previsíveis.}

\blocoNF
{NF008}
{Portabilidade}
{O sistema deve operar nos ambientes especificados com o mínimo de adaptações.}

% --------------------------------------------------
% 6 OUTROS REQUISITOS
% --------------------------------------------------
\section{Outros Requisitos}
Se necessário, incluir nesta seção requisitos complementares que não se enquadrem nas categorias anteriores, como requisitos legais, regulatórios, acadêmicos ou organizacionais.

\vfill

\end{document}